\documentclass[../../../../main.tex]{subfiles}
\begin{document}

$O$を原点とする座標平面上で考える。
座標平面上の2点$S(x_1, \, y_1)$，$T(x_2, \, y_2)$に対し，
点$S$が点$T$から十分離れているとは，
\[ |x_1-x_2|\geqq1 \hspace{1zw} \text{または} \hspace{1zw} |y_1-y_2|\geqq1 \]
が成り立つことと定義する。\\
　不等式
\[ 0 \leqq x \leqq 3, \hspace{1zw} 0 \leqq y \leqq 3 \]
が表す正方形の領域を$D$とし，その2つの頂点$A(3, \, 0)$，$B(3, \, 3)$を考える。
さらに，次の条件(i)，(ii)をともに満たす点$P$をとる。
\begin{enumerate}
  \item点$P$は領域$D$の点であり，かつ，放物線$y=x^2$上にある。
  \item点$P$は，3点$O, \, A, \, B$のいずれからも十分離れている。
\end{enumerate}
点$P$の$x$座標を$a$とする。
\begin{enumerate}
  \item$a$のとりうる値の範囲を求めよ。
  \item次の条件(iii)，(iv)をともに満たす点$Q$が存在しうる範囲の面積$f(a)$を求めよ。
\begin{enumerate}
  \item点$Q$は領域$D$の点である。
  \item点$Q$は，4点$O, \, A, \, B, \, P$のいずれからも十分離れている。
\end{enumerate}
  \item$a$は(1)で求めた範囲を動くとする。(2)の$f(a)$を最小にする$a$の値を求めよ。
\end{enumerate}

\end{document}
